\section{Discussion}\label{sec:discussion}

\subsection{Implications for knowledge distillation}

Our results demonstrate that a large fraction of the equational implication
structure can be captured by a small number of structural rules.  The
\emph{absorbing rule} alone (Theorem~\ref{thm:absorbing}) explains 46.8\% of
all true implications, and the \emph{general non-implication rule}
(Theorem~\ref{thm:general-nonimplication}) correctly classifies 22.2\% of all
pairs.  Together with the signature direction principle and variable count
heuristics, these rules achieve 88.9\% accuracy from purely syntactic features.

This is notable because the original proofs of many individual implications
in \cite{ETP2025} required sophisticated techniques: automated theorem provers,
SAT solvers, e-graph reasoning, and hand-crafted infinite constructions.  Yet
the \emph{aggregate} structure is largely predictable from simple invariants.
The gap between individual proof difficulty and aggregate predictability
suggests that equational implication, while undecidable in general
\cite{McNulty1976}, is statistically regular for low-order laws.

\subsection{Connections to the latent space approach}

Berlioz and Melli\`{e}s \cite{BM2026} showed that the Stone pairing --- the
probability that a random assignment satisfies an equation in a random finite
magma --- captures the implication structure in a low-dimensional latent space.
Our structural form classification is closely related: absorbing laws have
Stone pairing value zero on all nontrivial magmas (since they force $|M| = 1$),
while general-form laws tend to have higher satisfaction probabilities on random
magmas.

The variable count (Proposition~\ref{prop:nvar-monotone}) correlates strongly
with the first principal component of the Stone pairing matrix, which Berlioz
and Melli\`{e}s identify as the \emph{expected satisfaction probability}.
Laws with more variables are harder to satisfy, so they occupy lower positions
in the latent space --- and implications flow downward.

A natural question is whether the Stone pairing values themselves could be
encoded compactly to improve prediction.  Each law would require only a short
vector of probabilities (perhaps $3$--$5$ components from PCA), but encoding
$4{,}694$ such vectors may exceed the 10KB budget of the distillation challenge.

\subsection{Limitations}

Our approach has several important limitations.

\begin{enumerate}[label=(\roman*)]
\item \textbf{The hard middle:} Pairs with the same structural form, similar
  signatures, and equal variable counts remain difficult.  These constitute
  approximately 30\% of all pairs, and our accuracy on this subset is only
  55--65\%.  Improving accuracy on these ``hard middle'' cases likely requires
  algebraic reasoning beyond syntactic features.

\item \textbf{Computational steps and LLM reliability:} Steps~6 and~7 of the
  decision procedure (counterexample checking and rewriting) require symbolic
  computation that may not be reliably performed by a weak language model.
  The text-only cheatsheet version of our rules necessarily omits these steps,
  reducing accuracy from 88.9\% to approximately 82\%.

\item \textbf{Specificity to order $\le 4$:} Our analysis covers the $4{,}694$
  laws of order at most~$4$.  It is unclear whether the same structural rules
  generalize to higher-order laws.  The form classification extends naturally,
  but the calibrated probability estimates may differ substantially.

\item \textbf{Finite vs.\ infinite implications:} Our rules do not distinguish
  between $\models$ and $\models_{\mathrm{fin}}$.  By Austin's observation
  \cite{Austin1965}, these can differ: $E_i \models_{\mathrm{fin}} E_j$ does
  not imply $E_i \models E_j$ in general.  However, for the $4{,}694$ laws
  under consideration, the Equational Theories Project found extremely few
  discrepancies.
\end{enumerate}

\subsection{Open questions}

\begin{conjecture}\label{conj:95}
There exists a set of decision rules, expressible in at most $10$\,KB of text,
that achieves $\ge 95\%$ accuracy on the full implication matrix.
\end{conjecture}

The gap between our $88.9\%$ and the CNN's $99.7\%$ \cite{ETP2025} suggests
substantial room for improvement.  The CNN presumably learns features analogous
to our structural rules but also captures subtler patterns in term structure.
Distilling these learned features into interpretable rules is an open problem.

\begin{conjecture}\label{conj:equivalence}
Encoding the $1{,}415$ equivalence classes explicitly (by listing one
representative per class and the equivalence mappings) would improve accuracy
by at least $5$ percentage points.
\end{conjecture}

Such an encoding would allow the decision procedure to first reduce any law
to its canonical representative, collapsing many apparently distinct cases.
The storage cost is bounded by $1{,}415$ entries of approximately $3$~bytes
each (a class index), totaling about $4$\,kB.

We also pose the following theoretical question.

\begin{conjecture}\label{conj:undecidability}
For equational laws of unbounded order, no finite set of syntactic rules
achieves accuracy exceeding $80\%$ on the implication relation.
\end{conjecture}

This would follow if the implication relation for high-order laws becomes
sufficiently irregular that syntactic features lose their predictive power,
consistent with the general undecidability of the equational implication
problem \cite{McNulty1976}.
